%% Body only (paste between \begin{document} and \end{document}). Needs amsmath only.

\section{Experiment 1: The paper's method on 25 controlled models}

\textbf{Why.} In the paper, the models differ in architecture, tokenizer, size \emph{and} training
data at the same time, so any of these could explain the shape of the number line. Here only the
data differs.

\textbf{What we did.}
\begin{itemize}
\item \textbf{Models (DataDecide).} 25 models released by Ai2, all 1B parameters, with the same
  architecture, tokenizer and training length (100B tokens). Only the pretraining corpus (the
  ``recipe'') differs. Recipes form 6 families (Dolma, DCLM, Falcon, FineWeb, C4, DCLM+Dolma mixes);
  members of a family are filtered or modified versions of one corpus.
\item \textbf{$\beta$, exactly as in the paper.} Copy prompts \texttt{a=a,b=b,c=c,n=} with $n$ from four
  groups (1--39, 80--119, 980--1019, 9980--10019); only correctly answered prompts are kept. The
  last-token hidden state is projected on its first principal component (PC1), along which the
  numbers line up by size. $\beta$ is fitted from the 3 gaps between the four group positions:
  $\beta = 1$ is a logarithmic line, and a lower $\beta$ means large numbers are squeezed more. The
  layer is chosen by the paper's rule (the layer whose PC1 explains most variance and orders the
  numbers best) on one set of prompts, and $\beta$ is measured on fresh prompts.
\item \textbf{$\alpha$, as in the paper:} slope of $\log(\text{count of } n)$ against $\log n$ over
  $n = 1$--$10000$, counted on the exact 100B tokens each model was trained on. A less negative
  $\alpha$ means large numbers are relatively more common, so the paper predicts a positive
  $\alpha$--$\beta$ correlation.
\end{itemize}

\begin{table}[h]
\centering\small
\begin{tabular}{lccc}
\hline
Data & $n$ & $r$ ($p$) & $\rho$ ($p$) \\
\hline
Paper (corpus level) & 6 & 0.81 (0.061) & 0.60 (0.24) \\
DataDecide, all models & 25 & 0.56 (0.003) & 0.46 (0.022) \\
DataDecide, no Dolma & 19 & 0.22 (0.36) & 0.11 (0.65) \\
\hline
\end{tabular}
\caption{$\alpha$ vs $\beta$ in the paper and on DataDecide ($r$: Pearson, $\rho$: Spearman). All 25 models have a clear number line
(order of numbers along PC1: Spearman 0.90--0.96); $\beta$ ranges from 0.62 to 1.36.}
\label{tab:exp1}
\end{table}

\textbf{Result.} With architecture, tokenizer, size and training length held fixed, the predicted
positive relation appears and is statistically significant over 25 models. It is moderate, and most
of it comes from the Dolma family versus the rest (Experiment~3 addresses this).

\textbf{What it addresses.}
\begin{itemize}
\item \textbf{Confounding by model design} (H7EL, VRQv, a1nj): architecture, tokenizer, size and
  training length are identical across the 25 models.
\item \textbf{Too few corpora / weak statistics:} 25 models instead of 6 corpora; the correlation is
  significant.
\item \textbf{Partial correlation controlling for scale} (a1nj): not needed, since scale is constant.
\end{itemize}

\clearpage
\section{Experiment 2: Is $\beta$ reliable? (10 groups instead of 4)}

\textbf{Why.} The paper fits $\beta$ from only 3 gaps between 4 groups of numbers. Reviewers asked
whether this is stable and how good the fit is.

\textbf{What we did.} Same 25 models and prompts as Experiment~1, but 10 magnitude groups instead of
4 (evenly spaced on a log scale from 10 to 10000, i.e.\ 9 gaps), 100 correctly answered prompts per
group, and 3 random seeds. $\beta$ is fitted by one curve through all 3000 prompts,
$\text{PC1} = a + s\,\beta^{\log_{10} n}$. We then recompute $\alpha$ vs $\beta$ with this new $\beta$.

\textbf{Reliability.} Across seeds $\beta$ changes by only 0.015 (median s.d.), against a spread of
0.55--1.38 between models. The curve fits well (median $R^2 = 0.89$, range 0.81--0.97), and at the
same layer the new $\beta$ agrees with the paper's 3-gap $\beta$ at $r = 0.96$.

\begin{table}[h]
\centering\small
\begin{tabular}{lcc}
\hline
$\beta$ used & $r$ ($p$) & $\rho$ ($p$) \\
\hline
Old $\beta$ (3 gaps, Exp.~1) & 0.56 (0.003) & 0.46 (0.022) \\
New $\beta$, Exp.~1 layers & 0.63 (0.001) & 0.52 (0.008) \\
New $\beta$, layers re-chosen & 0.20 (0.33) & 0.19 (0.37) \\
New $\beta$, all at layer 8 & 0.75 ($<$0.001) & 0.71 ($<$0.001) \\
New $\beta$, all at one layer, 6--11 & 0.61--0.75 & 0.45--0.71 \\
New $\beta$, all at one layer, 2--4 & $-0.35$ to 0.00 & $-0.37$ to $-0.03$ \\
\hline
\end{tabular}
\caption{$\alpha$ vs $\beta$ before and after changing how $\beta$ is measured (25 models;
$r$: Pearson, $\rho$: Spearman). At every layer from 6 to 11, $r$ is significant after correcting
for testing 16 layers.}
\label{tab:exp2}
\end{table}

\textbf{Result.} The better-measured $\beta$ gives a slightly stronger relation at the same layers
(0.63 vs 0.56). When the layers are re-chosen it drops to 0.20: for 8 models two layers score within
8\% of each other, and the rule jumps to an early layer (2--4) with a very different $\beta$. When all
models are compared at the same middle layer, the relation is clear at every layer from 6 to 11. So
$\beta$ itself is measured reliably; the sensitive step is the choice of layer.

\textbf{What it addresses.}
\begin{itemize}
\item \textbf{$\beta$ from few points; stability and fit quality} (a1nj, point 3): stable across seeds,
  good fit, and consistent with the 3-gap $\beta$.
\item \textbf{Sensitivity to the layer:} near-tied layers can change the result; comparing all models at
  the same depth gives a clear relation at every middle layer.
\end{itemize}

\clearpage
\section{Experiment 3: $\alpha$ measured on the numbers $\beta$ uses}

\textbf{Why.} $\beta$ is measured only on four groups of numbers (1--39, 80--119, 980--1019,
9980--10019). The paper's $\alpha$ is a slope over \emph{all} numbers 1--10000, so 90\% of its points
are 4-digit numbers that $\beta$ never sees. The two measures therefore describe different numbers.

\textbf{What we did.} Same 25 models, same $\beta$ and layers as Experiment~1; no model dropped.
\begin{itemize}
\item $\alpha_{\mathrm{old}}$: the paper's slope over all numbers 1--10000.
\item $\alpha_{\mathrm{new}}$: for each of the 4 probed groups, take the median count of its numbers;
  $\alpha_{\mathrm{new}}$ is the slope of log(median count) against log(group centre), one point per
  group. (The median stops one very common number, such as 10000, from dominating a group.)
\end{itemize}
Example: $\alpha_{\mathrm{old}}$ says \texttt{fineweb-edu}'s numbers become rare faster than
\texttt{c4}'s ($-1.70$ vs $-1.56$), so it should be more compressed, yet its $\beta$ is higher (0.89
vs 0.62). On the probed numbers its counts actually fall more slowly, and $\alpha_{\mathrm{new}}$
orders the two correctly.

\begin{table}[h]
\centering\small
\begin{tabular}{lccc}
\hline
 & & $\alpha_{\mathrm{old}}$ & $\alpha_{\mathrm{new}}$ \\
Models used & $n$ & $r$ / $\rho$ & $r$ / $\rho$ \\
\hline
All 25 & 25 & .56$^{*}$ / .46$^{*}$ & \textbf{.82}$^{*}$ / \textbf{.77}$^{*}$ \\
No Dolma models & 19 & .22 / .11 & .68$^{*}$ / .66$^{*}$ \\
No code/math recipes & 17 & $-$.38 / $-$.41 & .44 / .52$^{*}$ \\
Only base corpora & 8 & .55 / .50 & .92$^{*}$ / .90$^{*}$ \\
Within families & 25 & .28 / .24 & .75$^{*}$ / .75$^{*}$ \\
Drop any one model ($r$) & 24 & .50--.61 & .79--.84 \\
\hline
\end{tabular}
\caption{$\alpha$ vs $\beta$ with both definitions of $\alpha$. $^{*}p<0.05$. ``Within families'':
each family's average removed first. ``Base corpora'': one recipe per original corpus, no
filtered variants. Dropping any one model: range of $r$ ($\rho$: .39--.52 old, .75--.80 new).}
\label{tab:exp3}
\end{table}

\textbf{Result.} Measured on the same numbers as $\beta$, $\alpha$ predicts $\beta$ much better
($r = 0.82$, $p < 10^{-6}$). Unlike $\alpha_{\mathrm{old}}$, it holds without Dolma, within families, and
on the independent base corpora, and no single model drives it. Because recipes in a family are
related, we also resampled whole families (bootstrap): the 95\% interval for $r$ is $[0.39, 0.90]$.

\textbf{Caveats.} This $\alpha$ was chosen after seeing that $\alpha_{\mathrm{old}}$ correlates weakly,
so it is exploratory and should be reported next to $\alpha_{\mathrm{old}}$ until tested on new
models. Without code/math recipes the relation is weaker ($r = 0.44$, $p = 0.08$), so data content
explains part of it.

\textbf{What it addresses.}
\begin{itemize}
\item \textbf{Sensitivity to individual corpora / leave-one-out} (VRQv Q2, H7EL): every leave-one-out
  $r$ is 0.79--0.84.
\item \textbf{Influence of Dolma:} $r = 0.68$ without the Dolma models.
\item \textbf{Analysis by family} (a1nj Q1): $r = 0.75$ within families.
\item \textbf{Uncertainty:} family-level bootstrap interval $[0.39, 0.90]$, excluding zero.
\end{itemize}
